% ============================================================
% Section 3: Data & Problem Setup
% ============================================================
\section{Data \& Problem Setup}
\label{sec:data}

\subsection{Data Sources}

We collect lecture materials from two Duke University courses: Deep Learning (5 lecture sets) and Machine Learning (5 lecture sets), totaling 10 lecture topics. For each topic, three complementary source types are gathered:

\begin{itemize}
    \item \textbf{Slide decks} (PPTX): Instructor-authored slides containing structured content, figures, and key definitions.
    \item \textbf{Auto-generated transcripts} (TXT): YouTube auto-captions (identified by the \texttt{[Auto-generated by YouTube]} header) from recorded lectures, providing spoken explanations and elaborations.
    \item \textbf{Student notes} (TXT): Handwritten or typed notes capturing a student's interpretation of the material.
\end{itemize}

These three sources are complementary: slides offer structured definitions and visual cues, transcripts capture verbal explanations and examples, and student notes reflect how learners organize and internalize the material. Combining all three gives summarization models a richer and more complete picture of each lecture.

For evaluation, we prepare one human-written reference summary per topic, yielding 10 reference summaries used to compute ROUGE-L and BERTScore.

\subsection{Data Processing Pipeline}
\label{sec:pipeline}

Our processing pipeline converts raw lecture materials into clean, model-ready text through four stages: cleaning, chunking, embedding, and topic assignment. Figure~\ref{fig:pipeline} illustrates the full flow.

\paragraph{Stage 1: Cleaning.}
Each source type undergoes tailored preprocessing:
\begin{itemize}
    \item \textbf{Transcripts}: We remove the YouTube auto-caption disclaimer header (e.g., \texttt{[Auto-generated by YouTube]}), join fragmented caption lines, collapse extra whitespace, and strip filler words (\textit{um}, \textit{uh}, \textit{like}, \textit{you know}, \textit{sort of}, etc.). Noise lines matching common patterns (e.g., greetings, audio checks, Q\&A markers, break announcements) are removed. The first and last 50 lines are additionally trimmed if they match these noise patterns.
    \item \textbf{Student notes}: Line endings are normalized and runs of three or more blank lines are collapsed to two.
    \item \textbf{Slides}: Text is extracted from every slide in the PPTX file. Slide titles are identified via PowerPoint's placeholder index. Decorative symbols (\texttt{\#}, $\blacktriangle$, $\diamondsuit$, $\bigstar$, $\bullet$) are stripped to produce plain text.
\end{itemize}

\paragraph{Stage 2: Chunking.}
Cleaned text is split into chunks of approximately 200 tokens with one overlapping sentence between consecutive chunks. Chunks smaller than 30 tokens are discarded. We use the BART tokenizer (\texttt{facebook/bart-large-cnn}) for token counting. The chunking strategy varies by source type:
\begin{itemize}
    \item \textbf{Transcripts}: Sentence-level sliding window using NLTK \texttt{sent\_tokenize}. Sentences are greedily accumulated until the next sentence would exceed the chunk size limit, at which point the chunk is flushed and the last overlapping sentence carries over to the next chunk.
    \item \textbf{Student notes}: Split at blank lines (paragraph boundaries). Short adjacent blocks are merged; blocks exceeding the chunk size fall back to the sentence sliding window.
    \item \textbf{Slides}: Each slide serves as a semantic unit. Short adjacent slides are merged. Oversized slides are sub-chunked via the sentence sliding window.
\end{itemize}

\paragraph{Stage 3: Embedding.}
All chunks are embedded using \texttt{all-MiniLM-L6-v2}~\citep{reimers2019sentence}, a 384-dimensional sentence transformer trained on over one billion sentence pairs. Embeddings are L2-normalized so that cosine similarity reduces to a dot product. Processing uses a batch size of 32.

\paragraph{Stage 4: Topic Assignment.}
We define one query string per lecture topic (e.g., ``deep learning'', ``supervised learning''). Each query and each chunk are embedded with MiniLM. Cosine similarity is computed between all query--chunk pairs, and each chunk is assigned to the topic with the highest similarity score. Chunks whose best score falls below a relevance threshold of 0.35 are discarded as off-topic noise. The surviving chunks are concatenated per topic to produce one combined source file per lecture, which serves as the input to all summarization models.

The final output of the pipeline is a set of 10 topic-level source files, each merging slide content, denoised transcript text, and student notes into a single document ready for summarization.


% ============================================================
% Section 5: Experiments and Results
% ============================================================
\section{Experiments and Results}
\label{sec:results}

We evaluate StudyLens through three experiments designed to isolate the effects of model architecture, summarization strategy, and output length. All experiments use the same 10 lecture topics and share a unified evaluation pipeline.

\subsection{Models}
\label{sec:models}

We evaluate seven models spanning three paradigms. Table~\ref{tab:models} summarizes each model's key properties.

\begin{table}[H]
\centering
\caption{Summary of all models evaluated in this study. Each encoder--decoder model was pre-trained on a general language modeling objective and then fine-tuned on a domain-specific summarization dataset.}
\label{tab:models}
\begin{tabular}{llcl}
\toprule
Model & Architecture & Context & Fine-tuning Data \\
\midrule
Naive (first-5) & Extractive heuristic & --- & --- \\
TF-IDF (top-15) & Classical ML & --- & --- \\
BART-CNN & Encoder--decoder & 1,024 & CNN/DailyMail news \\
Long-T5 & Encoder--decoder & 16,384 & BookSum (books) \\
BART-SAMSum & Encoder--decoder & 1,024 & SAMSum (dialogues) \\
LED-arXiv & Encoder--decoder & 16,384 & arXiv (academic papers) \\
Qwen2.5-7B & Instruction-tuned LLM & 128,000 & General instruction data \\
\bottomrule
\end{tabular}
\end{table}

\paragraph{Extractive Baselines.}
The \textbf{Naive} baseline selects the first 5 sentences from the input. \textbf{TF-IDF} ranks all input sentences by their TF-IDF importance scores and selects the top 15 highest-scoring sentences, returned in their original order.

\paragraph{Fine-Tuned Encoder--Decoder Models.}
We evaluate four pre-trained encoder--decoder models, each fine-tuned on a different summarization dataset:
\begin{enumerate}
    \item \textbf{BART-CNN}~\citep{lewis2020bart}: \texttt{facebook/bart-large-cnn} with a 1,024-token context window. Fine-tuned on CNN/DailyMail news articles, this is one of the most widely used summarization models.
    \item \textbf{Long-T5}~\citep{guo2022longt5}: \texttt{pszemraj/long-t5-tglobal-base-16384-book-summary} with a 16,384-token context. Fine-tuned on BookSum for long-document summarization, making it suited for lengthy lecture transcripts.
    \item \textbf{BART-SAMSum}: \texttt{philschmid/bart-large-cnn-samsum} with a 1,024-token context. Fine-tuned on the SAMSum dialogue summarization dataset. We included this model because lecture transcripts contain conversational elements that dialogue-trained models might handle better.
    \item \textbf{LED-arXiv}~\citep{beltagy2020longformer}: \texttt{allenai/led-large-16384-arxiv} with a 16,384-token context. Uses the Longformer attention mechanism and is fine-tuned on arXiv academic papers, making it a natural fit for technical lecture content.
\end{enumerate}

All four models use sentence-aware greedy chunking: input text is split into sentences using NLTK, and sentences are greedily packed into chunks until the model's context window is full. Each chunk is summarized independently. Single sentences exceeding the context limit are hard-truncated at the token level.

\paragraph{Instruction-Tuned LLM.}
\textbf{Qwen2.5-7B-Instruct}~\citep{qwen2025qwen25} has a 128K-token context window, which is large enough to process an entire lecture in a single pass without chunking. We load the model with 4-bit NF4 quantization via BitsAndBytes to fit within GPU memory. Generation uses greedy decoding (no sampling) with a slight repetition penalty of 1.1 for reproducibility. The output length is controlled by an \textit{output ratio} parameter, which sets the target number of generated tokens as a fraction of the input length. We sweep this ratio across six values: 1\%, 3\%, 6\%, 9\%, 15\%, and 30\%.

\subsection{Experimental Setup}

\paragraph{Experiment 1: Multi-Model Architecture Comparison.}
We compare all seven models listed above. The four encoder--decoder models use the \textit{concat} strategy (chunk summaries concatenated as the final output). Qwen is evaluated at its best-performing ratio (6\%). All models are evaluated on the same 10 lecture topics.

\paragraph{Experiment 2: Concat vs.\ Final Strategy.}
The encoder--decoder models support two summarization strategies:
\begin{itemize}
    \item \textit{Concat}: Each input chunk is summarized independently, and the chunk-level summaries are concatenated to form the output.
    \item \textit{Final}: A second summarization pass is run over the concatenated chunk summaries to produce a more coherent, re-synthesized output.
\end{itemize}
We evaluate both strategies for all four encoder--decoder models (BART-CNN, Long-T5, BART-SAMSum, LED-arXiv) plus Claude Sonnet~4 on a two-topic subset with full NLI evaluation.

\paragraph{Experiment 3: Output Ratio Sensitivity (Qwen).}
We sweep the Qwen output ratio across six values---1\%, 3\%, 6\%, 9\%, 15\%, and 30\%---to study how output verbosity affects summary quality. This experiment identifies the optimal ratio that balances conciseness with coverage.

\subsection{Evaluation Metrics}

We evaluate using three complementary metrics:
\begin{itemize}
    \item \textbf{ROUGE-L F1}~\citep{lin2004rouge}: Measures the longest common subsequence between the generated summary and the human reference, capturing surface-level overlap. Computed using the \texttt{rouge-score} library with stemming enabled.
    \item \textbf{BERTScore F1}~\citep{zhang2020bertscore}: Computes token-level cosine similarity using contextual embeddings from \texttt{roberta-large}, then aggregates into precision, recall, and F1. Baseline rescaling is applied. For long texts, both the reference and summary are chunked into segments of at most 450 words, aligned pairwise, and scored individually; the final score is the average across aligned pairs.
    \item \textbf{NLI Factual Consistency}: Uses a cross-encoder NLI model (\texttt{cross-encoder/nli-distilroberta-base}) to measure whether summary sentences are entailed by, neutral to, or contradicted by the source text. Each summary sentence is paired with the top-5 most similar source chunks (by Jaccard word overlap), and the maximum entailment probability is recorded. We report the entailment ratio (fraction of sentences with entailment $>0.5$), average entailment score, and contradiction ratio.
\end{itemize}

ROUGE-L and BERTScore measure alignment with the human reference. NLI measures faithfulness to the original source material.

\subsection{Results}

\subsubsection{Experiment 1: Architecture Comparison}

Table~\ref{tab:architecture} reports the average ROUGE-L F1 and BERTScore F1 across all 10 lecture topics for each model.

\begin{table}[H]
\centering
\caption{Average performance across 10 lecture topics (Experiment~1). Best results in bold. All encoder--decoder models use the \textit{concat} strategy.}
\label{tab:architecture}
\begin{tabular}{llccccc}
\toprule
Model & Type & Words & ROUGE-L P & ROUGE-L R & ROUGE-L F1 & BERTScore F1 \\
\midrule
Naive (first-5) & Extractive & 88 & 0.456 & 0.037 & 0.069 & $-$0.140 \\
TF-IDF (top-15) & Classical ML & 2,319 & 0.121 & 0.262 & 0.165 & $-$0.221 \\
\midrule
BART-CNN & Enc--Dec (1K) & 112 & 0.268 & 0.028 & 0.051 & $-$0.283 \\
Long-T5 & Enc--Dec (16K) & --- & --- & --- & --- & --- \\
BART-SAMSum & Enc--Dec (1K) & --- & --- & --- & --- & --- \\
LED-arXiv & Enc--Dec (16K) & 605 & 0.107 & 0.053 & 0.070 & $-$0.544 \\
\midrule
Qwen2.5-7B (6\%) & LLM (128K) & 606 & \textbf{0.348} & \textbf{0.179} & \textbf{0.233} & \textbf{$-$0.079} \\
\bottomrule
\end{tabular}
\end{table}

\noindent\textit{Note:} Long-T5 and BART-SAMSum were evaluated in the concat/final strategy experiment (Experiment~2, Table~\ref{tab:strategy}) on a two-topic subset rather than the full 10-topic benchmark.

Qwen achieves the highest ROUGE-L F1 (0.233) and BERTScore F1 ($-$0.079), outperforming all other models by a wide margin. The encoder--decoder models (BART-CNN and LED-arXiv) perform below even the naive baseline on BERTScore, reflecting the fragmentation caused by chunking long lecture inputs into small context windows. TF-IDF achieves moderate recall (0.262) by copying many sentences verbatim, but its precision is low (0.121), indicating that most copied sentences are not relevant to the reference summary. The naive baseline has high precision (0.456) but extremely low recall (0.037), producing very short outputs (88 words on average) that overlap with the reference only by chance.

\paragraph{ROUGE-L Comparison.}
Figure~\ref{fig:rouge_baselines} shows the ROUGE-L F1 for the four baseline models (without Qwen) evaluated across all 10 topics. TF-IDF leads among baselines, but all four score below 0.17.

\begin{figure}[H]
\centering
\includegraphics[width=0.85\linewidth]{figures/rouge_l_f1_baselines_from_csv.png}
\caption{ROUGE-L F1 across baseline architectures (10-topic average, without Qwen). TF-IDF leads among baselines but remains well below Qwen's 0.233.}
\label{fig:rouge_baselines}
\end{figure}

\paragraph{BERTScore Comparison.}
Figure~\ref{fig:bert_baselines} shows BERTScore F1 for the same four baselines. All models score below zero after baseline rescaling. LED-arXiv scores the worst ($-$0.544), likely because its academic-paper fine-tuning produces a writing style mismatched with lecture summary conventions.

\begin{figure}[H]
\centering
\includegraphics[width=0.85\linewidth]{figures/bertscore_f1_baselines_from_csv.png}
\caption{BERTScore F1 across baseline models (10-topic average). Encoder--decoder models show the weakest semantic alignment.}
\label{fig:bert_baselines}
\end{figure}

\paragraph{Combined ROUGE-L Comparison (All Models Including Qwen).}
Figure~\ref{fig:rouge_ratio} plots ROUGE-L F1 for all models side by side, with Qwen shown at each of its four evaluated ratios. This visualization makes clear the large gap between the LLM and all other approaches.

\begin{figure}[H]
\centering
\includegraphics[width=0.85\linewidth]{figures/rouge_l_f1_by_bmodel_ratio_20260311_112958.png}
\caption{ROUGE-L F1 across all models and Qwen output ratios. Qwen at 6\% and 9\% outperforms every other model by a large margin.}
\label{fig:rouge_ratio}
\end{figure}

\paragraph{Combined BERTScore Comparison (All Models Including Qwen).}
Figure~\ref{fig:bert_ratio} shows the same comparison for BERTScore F1. Qwen at 6--9\% is closest to zero (best), while baseline models are not visible at this scale due to their much lower scores.

\begin{figure}[H]
\centering
\includegraphics[width=0.85\linewidth]{figures/bertscore_f1_by_bmodel_ratio_20260311_112958.png}
\caption{BERTScore F1 across all models and Qwen output ratios. Qwen at 6--9\% achieves the best semantic similarity.}
\label{fig:bert_ratio}
\end{figure}

\subsubsection{Experiment 2: Concat vs.\ Final Strategy (All Four Encoder--Decoder Models)}

Table~\ref{tab:strategy} reports the results of both strategies for all four encoder--decoder models plus Claude Sonnet~4, evaluated on two lecture topics with full NLI scoring.

\begin{table}[H]
\centering
\caption{Concat vs.\ Final strategy comparison across all four encoder--decoder models and Sonnet~4 (2 topics, with NLI). Best ROUGE-L per model in bold.}
\label{tab:strategy}
\begin{tabular}{llccccc}
\toprule
Model & Strategy & Words & ROUGE-L F1 & BERTScore F1 & Entail.\ Ratio & Contra.\ Ratio \\
\midrule
BART-CNN & concat & 238 & \textbf{0.068} & $-$0.251 & 0.000 & 0.581 \\
BART-CNN & final & 44 & 0.010 & $-$0.194 & 0.000 & 0.700 \\
\midrule
Long-T5 & concat & 183 & \textbf{0.024} & $-$0.242 & 0.000 & 0.486 \\
Long-T5 & final & 96 & 0.018 & $-$0.251 & 0.000 & 0.786 \\
\midrule
BART-SAMSum & concat & 170 & \textbf{0.021} & $-$0.265 & 0.000 & 0.613 \\
BART-SAMSum & final & 116 & 0.013 & $-$0.253 & 0.000 & 0.583 \\
\midrule
LED-arXiv & concat & 752 & \textbf{0.058} & $-$0.580 & 0.000 & 0.394 \\
LED-arXiv & final & 252 & 0.031 & $-$0.508 & 0.000 & 0.513 \\
\midrule
Sonnet~4 & concat & 527 & \textbf{0.059} & \textbf{$-$0.032} & 0.042 & 0.857 \\
Sonnet~4 & final & 299 & 0.036 & $-$0.193 & 0.033 & 0.800 \\
\bottomrule
\end{tabular}
\end{table}

Several patterns stand out across all four encoder--decoder models:
\begin{itemize}
    \item The \textit{concat} strategy consistently outperforms \textit{final} on ROUGE-L F1 for every model. The second summarization pass in the \textit{final} strategy compresses already short chunk summaries further, losing information without gaining coherence.
    \item None of the encoder--decoder models achieve a nonzero entailment ratio. This means no summary sentence from BART-CNN, Long-T5, BART-SAMSum, or LED-arXiv is confidently supported by the source text.
    \item Contradiction ratios are high across the board (39--79\%), indicating that these fine-tuned models frequently generate statements that conflict with the input lecture content.
    \item BART-SAMSum, despite being fine-tuned on dialogues (and thus expected to handle lecture transcripts better), does not outperform standard BART-CNN. Its concat ROUGE-L (0.021) is actually lower than BART-CNN's (0.068).
    \item Long-T5, despite its larger 16K-token window, does not outperform BART-CNN on this subset. Its book-summary fine-tuning does not transfer well to lecture content.
    \item Claude Sonnet~4 achieves the best BERTScore ($-$0.032 concat) but also the highest contradiction ratio (85.7\%), showing that large LLMs produce fluent text that nonetheless introduces unsupported claims.
\end{itemize}

\subsubsection{Experiment 3: Output Ratio Sensitivity (Qwen)}

Table~\ref{tab:ratio} shows how varying the output ratio affects Qwen's summarization quality across the four ratios evaluated on all 10 topics.

\begin{table}[H]
\centering
\caption{Effect of output ratio on Qwen summarization quality (averaged over 10 topics). The 6\% ratio achieves the best performance.}
\label{tab:ratio}
\begin{tabular}{cccc}
\toprule
Ratio & Avg.\ Words & ROUGE-L F1 & BERTScore F1 \\
\midrule
6\%  & 606  & \textbf{0.233} & \textbf{$-$0.079} \\
9\%  & 610  & 0.233 & $-$0.079 \\
15\% & 1,032 & 0.212 & $-$0.083 \\
30\% & 2,606 & 0.157 & $-$0.126 \\
\bottomrule
\end{tabular}
\end{table}

The 6\% ratio produces the best results. Performance at 6\% and 9\% is nearly identical, with both ROUGE-L and BERTScore within 0.001 of each other. This suggests a quality plateau for concise outputs---the model has already captured the key content at 6\%, and allowing slightly more tokens (9\%) does not help. At 15\%, ROUGE-L drops by 9\% relative to the 6\% baseline, and at 30\% it drops by 33\%. The 30\% ratio produces summaries averaging 2,606 words---over four times longer than the 6\% output---yet scores worse on every metric because the extra length introduces redundancy and filler.

\paragraph{Why 6\% Works Best.}
The 6\% ratio strikes the right balance: it generates roughly 600 words, which is enough to cover all major lecture topics and subtopics without padding the output with repeated definitions or restated concepts. At this ratio, Qwen's ROUGE-L F1 (0.233) is 1.4$\times$ higher than TF-IDF (0.165) and 4.6$\times$ higher than BART-CNN (0.051), establishing a clear advantage for instruction-tuned LLMs on lecture summarization.

\paragraph{ROUGE-L by Ratio.}
Figure~\ref{fig:rouge_all_ratio} plots ROUGE-L F1 for all models and ratios side by side. The four Qwen bars dominate the chart, with a clear decline from 6\% to 30\%.

\begin{figure}[H]
\centering
\includegraphics[width=0.85\linewidth]{figures/rouge_l_f1_by_bmodel_ratio_20260311_112958.png}
\caption{ROUGE-L F1 across all models and Qwen output ratios (10-topic average). The 6\% and 9\% ratios perform nearly identically, while 30\% shows a clear drop.}
\label{fig:rouge_all_ratio}
\end{figure}

\paragraph{BERTScore by Ratio.}
Figure~\ref{fig:bert_all_ratio} shows the same comparison for BERTScore F1, with a sharper drop between 15\% and 30\%.

\begin{figure}[H]
\centering
\includegraphics[width=0.85\linewidth]{figures/bertscore_f1_by_bmodel_ratio_20260311_112958.png}
\caption{BERTScore F1 across all models and output ratios. Semantic similarity drops noticeably beyond 9\%.}
\label{fig:bert_all_ratio}
\end{figure}

\subsubsection{Faithfulness via NLI}

NLI evaluation was conducted on baseline models (10-topic average) and on all Qwen ratios.

\paragraph{Baseline Entailment.}
Figure~\ref{fig:nli_baselines} shows the average NLI entailment scores for the four baseline models. Extractive methods (Naive and TF-IDF) show higher entailment than the encoder--decoder models because they copy sentences directly from the source. TF-IDF achieves the highest entailment score among baselines, likely because its top-ranked sentences contain the most distinctive and verifiable claims. However, even verbatim copying does not guarantee high entailment---the NLI model evaluates sentence-level semantic relationships, and copied sentences may lack sufficient context when evaluated in isolation.

\begin{figure}[H]
\centering
\includegraphics[width=0.85\linewidth]{figures/nli_entailment_by_model_ratio.png}
\caption{Average NLI entailment scores across baseline models (10-topic average). Extractive methods (Naive, TF-IDF) achieve higher entailment than neural models (BART, LED-arXiv).}
\label{fig:nli_baselines}
\end{figure}

\paragraph{Qwen Entailment by Ratio.}
Figure~\ref{fig:nli_ratio} shows NLI entailment across all four Qwen output ratios. Qwen at 6\% and 9\% achieves the highest entailment scores (approximately 0.19), outperforming all baseline models. Entailment decreases slightly at 15\% and 30\%, consistent with the ROUGE-L and BERTScore trends---longer outputs introduce more sentences that are harder to ground in the source.

\begin{figure}[H]
\centering
\includegraphics[width=0.85\linewidth]{figures/nli_entailment_by_bmodel_ratio_20260311_112958.png}
\caption{NLI entailment across all models and Qwen output ratios. Qwen at 6--9\% achieves the highest entailment among all models.}
\label{fig:nli_ratio}
\end{figure}

\subsection{Summary of Findings}

Four key findings emerge from our experiments:

\begin{enumerate}
    \item \textbf{Instruction-tuned LLMs outperform all fine-tuned encoder--decoder models.} Qwen2.5-7B at 6\% achieves the best ROUGE-L F1 (0.233), BERTScore F1 ($-$0.079), and NLI entailment across all models tested---including four fine-tuned encoder--decoder models (BART-CNN, Long-T5, BART-SAMSum, LED-arXiv) and two extractive baselines.
    \item \textbf{Fine-tuning domain does not transfer to lectures.} Despite being fine-tuned on dialogues (BART-SAMSum), books (Long-T5), academic papers (LED-arXiv), and news (BART-CNN), none of these models produce good lecture summaries. The domain gap between their training data and lecture transcripts is too large.
    \item \textbf{Context window size matters more than model family.} Models that must chunk long lectures into small windows lose cross-slide coherence and produce fragmented outputs. Qwen's 128K-token window processes the full lecture in a single pass, avoiding this problem entirely.
    \item \textbf{The 6\% output ratio is optimal.} Increasing the output ratio beyond 9\% consistently degrades all metrics. The optimal operating point lies at 6\% of input length, producing summaries of roughly 600 words.
\end{enumerate}


% ============================================================
% Section 6: Error Analysis
% ============================================================
\section{Error Analysis}
\label{sec:errors}

We inspect model outputs across all three experiments and identify four recurring failure modes.

\paragraph{Chunk Fragmentation (All Four Encoder--Decoder Models).}
BART-CNN and BART-SAMSum process inputs in chunks of 1,024 tokens, while Long-T5 and LED-arXiv use 16,384-token chunks. When a lecture spans tens of thousands of tokens, even the larger-context models must split the input into multiple chunks. This causes three specific problems: (1)~information that spans multiple slides or transcript segments is split across chunks, so the model never sees the full argument; (2)~each chunk summary repeats setup context (e.g., restating the topic name), leading to redundant phrasing; and (3)~the concatenated output reads as a list of disconnected paragraphs rather than a coherent summary. The \textit{final} strategy was designed to address this by running a second summarization pass, but in practice it over-compresses already short chunk summaries---BART-CNN final produces only 44 words on average---losing content without improving coherence (Table~\ref{tab:strategy}).

\paragraph{Fine-Tuning Domain Mismatch.}
Each encoder--decoder model was fine-tuned on a specific domain, and none of these domains match lecture transcripts:
\begin{itemize}
    \item \textbf{BART-CNN} (news): Produces summaries in a journalistic style that does not capture technical definitions or mathematical formulas well.
    \item \textbf{Long-T5} (books): Generates narrative-style prose that lacks the structured format students expect from lecture summaries.
    \item \textbf{BART-SAMSum} (dialogues): Despite lectures being partially conversational, the model does not benefit from this---its ROUGE-L (0.021) is actually lower than BART-CNN's (0.068). Lecture ``dialogue'' is fundamentally different from chat conversations.
    \item \textbf{LED-arXiv} (academic papers): The closest domain match, yet it achieves the worst BERTScore ($-$0.580) among all models. Its summaries adopt an abstract-like writing style that does not align with how reference summaries are written.
\end{itemize}

\paragraph{Redundancy at High Output Ratios.}
When Qwen's output ratio is set to 30\%, summaries average 2,606 words and ROUGE-L drops to 0.157 (from 0.233 at 6\%). Manual inspection reveals that these long summaries repeat definitions and restate concepts in slightly different words. For example, a 30\% summary for the NLP~2 lecture restates the definition of word embeddings three times across different paragraphs. This repetition inflates word count without adding information, and it dilutes the density of relevant content that ROUGE and BERTScore can match against the reference.

\paragraph{High Contradiction Rates Across All Models.}
The NLI evaluation reveals that all models produce summaries with substantial contradiction ratios (Table~\ref{tab:strategy}). Among the four encoder--decoder models, contradiction ratios range from 39\% (LED-arXiv concat) to 79\% (Long-T5 final). None of the four encoder--decoder models achieve a nonzero entailment ratio---meaning that no generated sentence from BART-CNN, Long-T5, BART-SAMSum, or LED-arXiv is judged as clearly supported by the source. Even Claude Sonnet~4, which achieves the best BERTScore, has a contradiction ratio of 80--86\%. These contradictions arise for two reasons. First, the NLI model (DistilRoBERTa-base, 512-token limit) evaluates summary sentences against short source excerpts, and a summary sentence that synthesizes information from multiple parts of the lecture may appear to contradict any single excerpt. Second, abstractive models occasionally introduce unsupported claims.

\paragraph{Limitations of the Evaluation.}
Our NLI evaluation has its own limitations. The cross-encoder model has a 512-token context window, which restricts how much source material can be considered per summary sentence. The top-5 source chunk retrieval (by Jaccard overlap) may miss relevant passages that use different vocabulary. BERTScore values are negative after baseline rescaling, making absolute interpretation difficult---relative comparisons between models remain valid, but the scores do not indicate ``good'' or ``bad'' quality on an intuitive scale.


% ============================================================
% Section 7: Conclusion
% ============================================================
\section{Conclusion}
\label{sec:conclusion}

In this work, we presented StudyLens, a system for automatically summarizing university lectures by combining slide decks, transcripts, and student notes. We evaluated seven summarization approaches---a naive extractive baseline, TF-IDF sentence extraction, four fine-tuned encoder--decoder models (BART-CNN, Long-T5, BART-SAMSum, LED-arXiv), and an instruction-tuned LLM (Qwen2.5-7B-Instruct)---across 10 lecture topics from Deep Learning and Machine Learning courses, using ROUGE-L, BERTScore, and NLI factual consistency as evaluation metrics.

Our experiments lead to four main conclusions:
\begin{enumerate}
    \item \textbf{LLMs outperform fine-tuned models.} Qwen2.5-7B at the 6\% output ratio consistently achieves the best results on all metrics. Its ROUGE-L F1 of 0.233 is 1.4$\times$ higher than TF-IDF (0.165), 3.3$\times$ higher than LED-arXiv (0.070), and 4.6$\times$ higher than BART-CNN (0.051).
    \item \textbf{Fine-tuning domain matters, but none of the available domains match lectures.} Despite trying models fine-tuned on news (BART-CNN), books (Long-T5), dialogues (BART-SAMSum), and academic papers (LED-arXiv), none produced good lecture summaries. The domain gap between training data and lecture transcripts is too large for fine-tuning alone to bridge.
    \item \textbf{Context window is more important than model architecture.} All four encoder--decoder models struggle because they must chunk long lectures. Qwen's 128K-token window avoids this entirely, processing the full lecture in a single pass.
    \item \textbf{The 6\% output ratio is optimal for Qwen.} Sweeping across six ratios (1\%--30\%), we find that 6\% produces roughly 600-word summaries that maximize both reference alignment and semantic quality. Beyond 9\%, quality degrades steadily.
\end{enumerate}

We also identified areas where all models fall short. Factual consistency remains a challenge: contradiction ratios are high even for the best-performing models, and all four encoder--decoder models achieve zero entailment. This gap between fluency and faithfulness suggests that lecture summarization would benefit from grounding mechanisms such as retrieval-augmented generation or constrained decoding.

Our study has several limitations. The evaluation is conducted on 10 lecture topics from a single institution, which limits generalizability. The human reference summaries were written by one annotator per topic, introducing potential bias. NLI evaluation uses a relatively small cross-encoder model with a 512-token window, which may undercount entailment for summaries that synthesize information from multiple source passages. Long-T5 and BART-SAMSum were evaluated on a two-topic subset rather than the full 10-topic benchmark due to computational constraints.

Based on our findings, we offer three practical recommendations for lecture summarization systems:
\begin{enumerate}
    \item Use full-context LLMs that can process the entire lecture in a single pass, avoiding the chunking artifacts that degrade encoder--decoder models.
    \item Set the output ratio between 6\% and 9\% of input length to balance coverage and precision.
    \item Integrate retrieval-grounded verification to reduce the rate of unsupported claims in generated summaries.
\end{enumerate}
